\section{Dynamic and 4D Reconstruction}
\label{sec:dynamic}

Dynamic-scene Event-GS is the most methodologically active category. The common substrate is a Deformable 3DGS (deformation MLP + HexPlane grid, or 4D primitives) that predicts per-Gaussian deltas at query time $t$. Events provide the high-temporal-resolution motion supervision that RGB cannot, but raise three challenges: (i) how to query arbitrary timestamps given sparse poses; (ii) how to constrain photometry without absolute intensity; (iii) how to correct depth under monocular sparse views.

\subsection{Event-Boosted Deformable 3D Gaussians (ICCV 2025)}
\label{sec:dynamic:boosted}

Event-Boosted~\cite{xu2025eventboosted} is the first event-camera + deformable-3DGS combination for dynamic scenes. It makes two contributions. (1) \emph{GS-Threshold Joint Modeling (GTJM)}: observing that event threshold modeling is critical to reconstruction quality, it jointly optimizes the Gaussian field and the event threshold, with each informing the other---eliminating the ``purple haze'' artifact caused by a misspecified fixed threshold. (2) \emph{Dynamic-Static Decomposition (DSD)}: dynamic areas are identified by the inability of static Gaussians to represent motion, then a buffer-based soft decomposition separates dynamic and static Gaussians, accelerating rendering and focusing deformation capacity on dynamic regions. It contributes the first event-inclusive 4D benchmark with synthetic and real-world dynamic scenes.

\textbf{Core point:} the event threshold is a learnable, jointly-optimized parameter, not a fixed constant---an early step toward differentiable event-camera physics.

\subsection{FastEventDGS (CVPR 2026)}
\label{sec:dynamic:fast}

FastEventDGS~\cite{dai2026fasteventdgs} is the first method to train Deformable 3DGS for dynamic scenes using \emph{only a single monocular event camera}, with no auxiliary RGB sensor. Its four design pillars are:
\begin{itemize}
    \item \textbf{Continuous trajectory parameterization:} cubic B-splines fit sparse poses into a continuous trajectory, enabling event supervision at any timestamp $\tau$---essential because event frequency vastly exceeds pose sampling rate.
    \item \textbf{Dual event generation model:} an Event Single Integral (ESI) loss provides global photometric consistency, while a linear event-flow loss $\Delta I(\mathbf{x})\approx-\nabla I\cdot\mathbf{v}\,\Delta\tau$ provides local geometric constraints via camera flow + Gaussian flow.
    \item \textbf{Local patch motion loss:} projects sampled Gaussian centers to the image plane and enforces that consecutive event integrals on patches around them are consistent, pinning Gaussian trajectories to event edges and suppressing background jitter.
    \item \textbf{Refinement stage:} uses the VGGT feed-forward model as a depth expert with a Scale-Invariant Logarithmic (SiLog) loss, plus non-event-region and spatial-smoothness regularizers.
\end{itemize}
It improves PSNR from $\sim$16\,dB to 22--24\,dB on synthetic and real fast-motion datasets.

\textbf{Core point:} events should be treated as explicit supervision for the motion field, not merely brightness differences; the combination of global integration + local flow + motion-edge pinning resolves the monocular under-constraint.

\subsection{ERF-GS (CVMJ 2026)}
\label{sec:dynamic:erfgs}

ERF-GS~\cite{bai2026erfgs} (the focus of this project's codebase) reconstructs fast motion from \emph{disjoint} multi-view Event-RGB viewpoints: some cameras supply only RGB, others only events, never both. Its three contributions are:
\begin{itemize}
    \item \textbf{linlog + STE event loss:} predicted events are the linlog-luma difference between two deformed-frame renders (Eq.~\eqref{eq:linlog}); the loss uses a straight-through estimator (STE) to quantize predictions to integer polarity counts while preserving gradients, normalized by $\max(|e_{gt}|,1)$.
    \item \textbf{Linear-motion + constant-color regularization:} between event-interval endpoints, the deformed position is required to follow a linear interpolation and the color to remain constant, encoding a constant-velocity prior.
    \item \textbf{Event-driven densification:} event pixels are back-projected at multiple candidate depths across views, filtered by multi-view consistency, selected by deformation-aware isolation, and initialized via motion-warped KNN. This spawns new Gaussians in fast-moving regions that RGB-gradient densification misses.
\end{itemize}
The disjoint-viewpoint setup forces cross-modal fusion and is the paper's defining experimental contribution.

\textbf{Core point:} events drive not only the loss but also the \emph{densification}---a capability unique to explicit GS representations.

\subsection{PEGS (arXiv 2025)}
\label{sec:dynamic:pegs}

PEGS~\cite{xu2025pegs} brings \emph{physics priors} to event-driven GS for large spatiotemporal-scale rigid-body motion. It layers three supervisions: (i) acceleration-consistency constraints from Newtonian dynamics, (ii) high-temporal-resolution event-flow guidance, and (iii) Kalman regularization fusing multi-source observations. A motion-aware simulated-annealing strategy schedules the constraints. It contributes the first paired RGB-Event dataset of natural fast rigid-body motion.

\textbf{Core point:} physics priors (acceleration consistency) and event supervision are complementary---events give high-rate positions, physics gives second-order smoothness---and their combination is, to our knowledge, unique to PEGS within the Event-GS family.

\subsection{Comparison of dynamic methods}

\begin{table}[t]
\centering
\caption{Comparison of dynamic/4D Event-GS methods.}
\label{tab:dynamic}
\small
\begin{tabular}{@{}p{2.3cm}p{1.1cm}p{1.1cm}p{1.6cm}p{1.0cm}@{}}
\toprule
Method & Input & Pose & Event model & Physics \\
\midrule
Event-Boosted~\cite{xu2025eventboosted} & E+R & Known & Learnable thresh. & -- \\
FastEventDGS~\cite{dai2026fasteventdgs} & E & Sparse & ESI+flow & -- \\
ERF-GS~\cite{bai2026erfgs} & E+R & Known & linlog+STE & Linear \\
PEGS~\cite{xu2025pegs} & E+R & Known & Event flow & Accel. \\
\bottomrule
\end{tabular}
\end{table}

Table~\ref{tab:dynamic} shows the diversity: only FastEventDGS is event-only; only PEGS uses physics; only ERF-GS uses event-driven densification; only Event-Boosted jointly models the threshold. This four-way design space---input modality, pose regime, event model, prior type---maps the open frontier of dynamic Event-GS.
